% from ega4-9-fr.tex
\oldpage[IV-3]{90}

point générique de $(Z_j)_\eta$, donc
$\dim((X_i\cap Z_j)_\eta)<\dim((Z_j)_\eta)$~; par suite, la dimension commune
des composantes irréductibles de $(X_i)_s\cap(Z_j)_s$ est, pour tout $s\in S$,
strictement inférieure à la dimension commune des composantes irréductibles de
$(Z_j)_s$, ce qui prouve qu'aucune des composantes irréductibles de $(Z_j)_s$
n'est contenue dans une composante irréductible de $(X_i)_s$ pour $s\in S$.
C.Q.F.D.

\begin{proposition}[9.9.2]
\label{IV.9.9.2-fr}
Soient $f:X\to S$ un morphisme localement de présentation finie, $\mathcal F$
un $\mathcal O_X$-Module quasi-cohérent de présentation finie, $\Phi$ une
partie finie de $\mathbf Z\cup\{\pm\infty\}$. Les parties suivantes de $X$
sont localement constructibles~:
\begin{enumerate}
\item[(i)] L'ensemble des points $x\in X$ tels que $\mathcal F_{f(x)}$ soit
équidimensionnel au point $x$ (5.1.12).
\item[(ii)] L'ensemble des $x\in X$ tels que les longueurs géométriques de
$\mathcal F_{f(x)}$ relatives à $k(f(x))$, aux points génériques des composantes
irréductibles de $\operatorname{Supp}(\mathcal F_{f(x)})$ contenant $x$
(4.7.5), appartiennent à $\Phi$.
\item[(iii)] L'ensemble des $x\in X$ tels que les dimensions des cycles
premiers associés à $\mathcal F_{f(x)}$ et contenant $x$ appartiennent à
$\Phi$.
\item[(iv)] L'ensemble des $x\in X$ tels que $\mathcal F_{f(x)}$ soit
géométriquement réduit au point $x$ (4.6.17).
\item[(v)] L'ensemble des $x\in X$ tels que $\mathcal F_{f(x)}$ soit
géométriquement intègre au point $x$ (4.6.22).
\item[(vi)] L'ensemble des $x\in X$ tels que
$\operatorname{dim.proj}((\mathcal F_{f(x)})_x)\in\Phi$.
\item[(vii)] L'ensemble des $x\in X$ tels que
$\operatorname{coprof}((\mathcal F_{f(x)})_x)\in\Phi$.
\item[(viii)] L'ensemble des $x\in X$ tels que $\mathcal F_{f(x)}$ possède la
propriété $(S_k)$ au point $x$ (5.7.2).
\item[(ix)] L'ensemble des $x\in X$ tels que $\mathcal F_{f(x)}$ soit un
$\mathcal O_{X_{f(x)}}$-Module de Cohen-Macaulay au point $x$ (5.7.1).
\end{enumerate}
\end{proposition}

(i) Le support $Z$ de $\mathcal F$ est localement constructible et fermé dans
$X$ (8.9.1), et en considérant un sous-préschéma de $X$ ayant $Z$ pour espace
sous-jacent et qui soit de présentation finie sur $S$ (8.9.1), on voit que
l'assertion (i) est un cas particulier de (9.9.1, (iii)).

(ii) Toutes les propriétés envisagées sont locales sur $X$, et on se bornera
donc au cas où $X=\operatorname{Spec}(B)$ et $S=\operatorname{Spec}(A)$ sont
affines et $f$ un morphisme de présentation finie. On garde les notations du
début de la démonstration de (9.9.1), et on suppose en outre $A_0$ choisi de
sorte qu'il existe un $\mathcal O_{X_0}$-Module cohérent $\mathcal F_0$ tel que
$\mathcal F$ soit isomorphe à $\mathcal F_0\otimes_{A_0}A$. Alors ((4.2.7) et
(4.7.9)) l'ensemble des longueurs géométriques de
$(\mathcal F_0)_{f(x_0)}$ aux points génériques des composantes irréductibles de
$\operatorname{Supp}((\mathcal F_0)_{f(x_0)})$ qui contiennent $x_0$ est le
même que l'ensemble analogue pour $x$ et $\mathcal F_{f(x)}$~; autrement dit,
si $E$ (resp. $E_0$) est l'ensemble des $x\in X$ (resp. des $x_0\in X_0$)
vérifiant la condition (ii) de l'énoncé, on a $E=h^{-1}(E_0)$, et en vertu de
(1.8.2), on voit qu'on peut se borner à considérer le cas où $A$ est
\emph{noethérien}. Comme dans la démonstration de (9.9.1), on voit qu'on est
ramené à montrer que, pour tout $x\in X$, il existe un voisinage $V$ de $x$
dans $\overline{\{x\}}$ tel que, pour tout $x'\in V$, l'ensemble des longueurs
géométriques de $\mathcal F_{f(x')}$ aux points génériques des composantes
irréductibles de son support contenant $x'$ soit \emph{le même}. En outre, si
$S'$ est le sous-préschéma réduit de $S$ ayant $\overline{\{f(x)\}}$ pour espace
sous-jacent, et si $X'=f^{-1}(S')$, les fibres de $X$ et de $X'$ aux points de
$S'$ sont les mêmes,

